\documentclass[12pt]{article}
\usepackage[utf8]{inputenc}
\usepackage[T1]{fontenc}
\usepackage{hyperref}
\title{Bulo Stieng Language Resources for K-12 Curriculum Development}
\author{Compiled by Claude Sonnet 5.5 (claude-sonnet-5-5) for Andr\'as Kornai}
\date{}
\begin{document}
\maketitle

\section{Introduction}
Bulo Stieng (ISO 639-3 sti) is an Austroasiatic language spoken in Vietnam and adjacent parts of Cambodia. Wikipedia gives about 90,000 speakers for Stieng overall, and the Endangered Languages Project rates it as threatened. Little material was found in search results. Bunong (Mnong), which has a published dictionary, is a different language and should not be confused with Bulo Stieng.

\section{Primary Educational Resources}
No primers, textbooks or school materials for Bulo Stieng were found in search results.

\section{Linguistic and Reference Materials}
\begin{itemize}
\item Stieng language, Wikipedia. \url{https://en.wikipedia.org/wiki/Stieng_language}
\item Bulo Stieng, Endangered Languages Project bibliography. It lists the Haupers' work, including a dictionary of the Bulo dialect compiled between 1960 and 1975. I could not confirm whether it is in print or online. \url{https://endangeredlanguages.com/lang/5422/bibliography?hl=zh-cn}
\item Bulo Stieng page, Verbix Endangered Languages. \url{https://docs.verbix.com/EndangeredLanguages/BuloStieng}
\item Academic document hosted on HAL (Lyon 2). I did not verify its content. \url{https://hal.univ-lyon2.fr/hal-03025567/document}
\end{itemize}

\section{Religious and Community Resources}
\begin{itemize}
\item Joshua Project, Bulo Stieng in Vietnam (people-group profile). \url{https://www.joshuaproject.net/index.php/people_groups/15088/VM}
\item Joshua Project, Bulo Stieng in Cambodia. \url{https://joshuaproject.net/index.php/people_groups/15088/CB}
\end{itemize}

\section{Digital and Audio Resources}
No audio or corpus resources for Bulo Stieng were found. A Webonary Bunong-Khmer-English dictionary appeared in results, but it covers Bunong, a related but different language, so it is not listed as a Bulo Stieng resource.

\section{Curriculum Development Context}
No evidence of schooling in Bulo Stieng was found. A curriculum would have to start from the Haupers' dictionary material and from work with community members.

\end{document}